Fixation of the Residual Intensity  
Derived from the 6.5-Winding Trace

1. The 6.5-Winding Trace as Primary Source

The residual intensity is not postulated; it is read off a concrete geometric trace that originates in the maximal winding \(\theta_{\max}=6.5\pi\) (\(w=3.25\)).

On the unit-circle / complex-plane representation the terminal ray of the asymptotic direction

\[
\lim_{t\to 0^+}\arg\bigl(t(1+2\pi i)\bigr)=\arctan(2\pi)\approx 1.412965136507
\]

is drawn as a straight radial line from the origin. Its Euclidean length, when the argument is treated as a radial coordinate, is exactly \(1.412965136507\).

Independently, the thinnest residual near-equilateral triangle generated by the practical step \(7\psi\) possesses altitude

\[
h_{\mathrm{res}}=1.490969476641.
\]

When both quantities are rendered as distances from the origin in the same plane, the linear gap between them is

\[
1.490969476641-1.412965136507=0.078004340134.
\]

The corresponding ratio is

\[
\frac{h_{\mathrm{res}}}{\arctan(2\pi)}\approx 1.055206.
\]

This difference \(0.078004340134\) is the traced connection: the extra radial length that separates the terminal-ray argument of the 6.5-winding from the altitude of the residual triangle. It is the first quantitative signature that the residual geometry has detached itself from the pure terminal ray while remaining anchored to the same winding family.

2. Propagation of the Thinnest Near-Equilateral Triangle

The same residual generator that produces the altitude \(h_{\mathrm{res}}\) produces the full metric of the triangle \(T\):

\[
\begin{align*}
s_{\mathrm{thin}}&=1.721623257385,\\
s_{\mathrm{long}}&=1.737195272475,\\
\operatorname{Area}(T)&=1.2988990741.
\end{align*}
\]

Because the generator \(7\psi=\pi/3+\delta\) continues to act, congruent copies of this triangle repeat at every scale of the unit-circle hierarchy. The triangle is therefore not a local accident; it is the elementary residual cell that tiles, under residual radial folding, every subsequent densification.

3. Definition of the Dimensionless Residual Intensity

Inside the unit disk the complementary region that remains after \(T\) is excised has area \(1.8426935795\). Partitioned into three equal residual slices this yields

\[
A_{\mathrm{slice}}=0.6142311932.
\]

Their ratio to the shoelace area of the triangle defines the dimensionless residual intensity

\[
\rho=\frac{A_{\mathrm{slice}}}{\operatorname{Area}(T)}\approx 0.472886.
\]

Thus \(\rho\) is fixed by the same residual cell whose altitude already stood in the traced gap of \(0.078\) relative to the 6.5-winding terminal ray. The intensity inherits its origin from that trace.

4. Visibility as Edge Curvature

Each side of the residual triangles is originally a straight chord. The Core Relations convert that chord into a curved edge by a normal displacement proportional to the intensity just obtained:

\[
\gamma_\rho(t)=\gamma_0(t)+c\cdot\rho\cdot 4t(1-t)\,n,\qquad c=0.18.
\]

The resulting bulge \(\approx 0.08512\) is the direct geometric transcription of the residual area that remains after the shoelace triangle is removed from the unit disk. Consequently the curved sides of the Order-6 figure are the first concrete appearance of the residual intensity—itself descended from the 6.5-winding altitude-to-terminal-ray gap—as continuous edge curvature.

5. Closure of the Corrected Fixation

The residual intensity \(\rho\) is fixed by a continuous chain that begins with the 6.5-winding terminal ray, passes through the measured altitude gap \(0.078004340134\), determines the thinnest near-equilateral triangle, extracts its residual area slices, and forms the dimensionless ratio \(\rho\). When this intensity is applied as edge curvature, the Order-6 radial fold becomes the geometric locus at which the original winding trace is finally visible as both a curved boundary and an inner hexagonal measure. All higher densifications merely replicate that same intensity at finer combinatorial resolution.
